\section{Monotone boxes, direction obstructions, and robust rectangle folds}\label{sec:monotone}

This section records a second infinite-family result.  It concerns a
precisely restricted version of possible play, and is therefore separate
from the still-unresolved unrestricted rectangle construction discussed
in Section~\ref{sec:recent}.  Fix one direction on each nontrivial axis of
a finite box and allow only those directions throughout the play.  By
reflecting axes, call them the negative directions.  On a rectangle these
are Up and Left.  Spawns remain unrestricted: after every changed slide,
a 2 or 4 may appear in any empty cell.

Let the box have side lengths $n_1,\ldots,n_d$ and at least two cells.  Put
\[
 c=\prod_{i=1}^d n_i,\qquad K_i=2^{n_i}-1,\qquad
 K=\prod_{i=1}^d K_i,
\]
and for a cell $x=(x_1,\ldots,x_d)$, $0\le x_i<n_i$, define
\[
 A(x)=2+\sum_{i=1}^d(n_i-1-x_i),\qquad
 D(x)=2^{A(x)}.
\]
Thus every step away from the chosen corner halves $D$.  We continue to
write $\Phichain(B)=\sum_{2^k\in B}(k-1)2^k$ for the score potential of a
position.

\begin{theorem}[Exact monotone extrema]\label{thm:monotone}
Under possible play using only the negative direction on each axis,
\begin{align*}
 \max\text{ tile}&=2^{\,2+\sum_i(n_i-1)},\\
 \max\mass&=4K,\\
 \max\mathrm{score}&=\Phichain(D)-4c,\\
 \max\text{ rounds}&=2K-c-2.
\end{align*}
The full geometric array $D$ is reachable from an ordinary two-2 opening
with exactly $c$ spawned 4s; this play simultaneously attains the score
and round ceilings.  It is also reachable from an ordinary two-4 opening
using only 4-spawns in exactly $K-2$ rounds, which is the fewest rounds in
which the complete array $D$ can be reached.

If
\[
 J_i=(n_i-2)2^{n_i}+2,
\]
then
\[
 \Phichain(D)=4\left[K+\sum_iJ_i\prod_{j\ne i}K_j\right].
\]
The minimum $K-2$ concerns the full array $D$, not necessarily the first
appearance of its largest tile.
\end{theorem}

For a rectangle of height $h$ and width $w$ this gives
\[
 \max\text{ tile}=2^{h+w},\qquad
 \max\text{ rounds}=2(2^h-1)(2^w-1)-hw-2,
\]
and
\[
 \max\mathrm{score}=4\big[(h+w-3)2^{h+w}+(3-h)2^h
 +(3-w)2^w-3-hw\big].
\]
For example, the exact Up/Left extrema on $4\times4$ are tile 256,
score 4916, and 432 rounds.  The saturated array is
\[
\begin{pmatrix}
256&128&64&32\\
128&64&32&16\\
64&32&16&8\\
32&16&8&4
\end{pmatrix},
\]
of mass 900; its fastest construction has 223 rounds.  This is a theorem
about the restricted game, not a claim that the unrestricted $4\times4$
score is 4916.

\subsection{A pointwise directional bound}

For a tile of rank $k$ at $x$, set
\[
 p=k+\sum_i x_i.
\]
A nonmerging negative-direction movement cannot increase $p$.  Suppose a
merge along axis $i$ has equal-rank inputs at coordinates $u<v$ on that
axis and output coordinate $q$.  Since $q\le u$,
\[
 (k+1)+q\le k+v,
\]
so the output potential is no larger than that of one input.  Because a
tile merges at most once in a slide, the same statement holds along every
ancestry.

\begin{lemma}[Ancestry-sensitive pointwise bound]\label{lem:monopoint}
Let $h\in\{0,1\}$ indicate whether the merge tree of a tile at $x$ has a
spawned 4 among its leaves.  Its rank satisfies
\[
 k\le A(x)-1+h.
\]
In particular every cell satisfies $B(x)\le D(x)$.
\end{lemma}
\begin{proof}
A spawned 2 has initial potential at most
$1+\sum_i(n_i-1)$ and a spawned 4 at most
$2+\sum_i(n_i-1)$.  The potential cannot increase along a lineage under
the permitted slides.  Subtracting $\sum_i x_i$ gives the claim.
\end{proof}

The mass bound of Theorem~\ref{thm:monotone} follows immediately and
uniquely identifies the saturated spatial array, not merely its multiset.
For score and length, let $h_x$ be the indicator of a 4-leaf in the final
tile at $x$.  Final-tile ancestries are disjoint, so
$\nfour\ge\sum_xh_x$.  At a cell with $a=A(x)\ge2$, Lemma~\ref{lem:monopoint}
and the score identity give the local upper bound
\[
 \max_{k\le a-1+h}\big((k-1)2^k-4h\big)=(a-1)2^a-4.
\]
For $a>2$ the unique maximizing choice is $(k,h)=(a,1)$; for $a=2$,
$(2,1)$ and $(1,0)$ tie at zero.  Likewise
\[
 \max_{k\le a-1+h}\big(2^{k-1}-h\big)=2^{a-1}-1.
\]
Summing proves the score and round ceilings.  Equality has the same
last-cell tie familiar from Theorem~\ref{thm:score}: the far-corner 4 may
be replaced by a final 2, saving one 4-birth without changing score or
round count.

\subsection{A search-free construction and its invariant}

The attainment proof is recursive.  Its essential one-dimensional
counter is useful independently.

\begin{lemma}[Two scalar counters]\label{lem:monocounter}
On a line of $n\ge2$ cells, keep sliding toward one end and place each
new spawn at the far end.
\begin{enumerate}[nosep,label=(\roman*)]
\item Starting from a single 2 and always spawning 2 terminates at
$(2^n,2^{n-1},\ldots,2)$.
\item Starting from a single 2, spawn 4 precisely when the afterstate has
exactly $n-1$ nonzero entries, all distinct and all at least 4; otherwise
spawn 2.  The terminal line is $(2^{n+1},2^n,\ldots,4)$, and exactly $n$
spawned 4s are used.
\item Starting from a single 4 and always spawning 4 also terminates at
$(2^{n+1},2^n,\ldots,4)$.
\end{enumerate}
\end{lemma}
\begin{proof}
Ignoring zeros, entries remain nonincreasing.  After every spawn the far
cell is occupied, so the sole direction is blocked exactly when the line
is full and strictly decreasing.  The pointwise bound limits the largest
rank; therefore the pure-2 and pure-4 terminal lines are forced.

For the adaptive rule, if a 2-spawn would complete a blocked line, the
afterstate has precisely the property that triggers a 4.  If its smallest
entry is 4, that 4 creates another merge; if its smallest entry is at
least 8, the $n-1$ entries are exactly $2^{n+1},\ldots,8$, and the spawned
4 completes the claimed terminal line.  Each time the rule chooses 4, the
afterstate is an $(n-1)$-element subset of
$\{4,8,\ldots,2^{n+1}\}$; different choices occur at different masses, so
there are at most $n$ such events.  Conversely the $n$ terminal tiles all
saturate their pointwise bound, hence each needs a disjoint 4-leaf by
Lemma~\ref{lem:monopoint}.  Thus exactly $n$ fours are necessary and used.
\end{proof}

For an inner box with coordinates $y$, define its unit profile
\[
 E(y)=\prod_i2^{n_i-1-y_i}.
\]
A completed slice has the form $vE$.  Every inner-axis line in such a
slice is full with pairwise distinct entries, so every inner-axis
construction move leaves the completed slice unchanged.  Across an outer
axis, slices $v_0E,\ldots,v_{m-1}E$ slide and merge in exact lockstep:
each physical line sees the same coefficient list multiplied by one fixed
entry of $E$.  Hence an outer slide is exactly a scalar 2048 slide on
$v_0,\ldots,v_{m-1}$.

\begin{lemma}[Recursive slice construction]\label{lem:monoslices}
For every nontrivial box, from a single 2 at the far corner the profile
$2E$ is constructible with no 4-spawns and $4E$ is constructible with
exactly one spawned 4 per cell.  From a single 4, $4E$ is constructible
using only 4-spawns.  All moves use the chosen negative directions and
all physical spawns occur at the far corner.
\end{lemma}
\begin{proof}
Induct on the number of nontrivial axes.  Lemma~\ref{lem:monocounter} is
the one-dimensional base.  For the induction step, first construct the
required inner profile in the far outer slice.  Treat each completed
slice $vE$ as the scalar tile $v$ and run the corresponding scalar counter
along the outer axis.  Each outer slide empties the far slice.  Its
mandatory physical spawn starts the recursively constructed replacement
slice there: a pure inner construction when the virtual counter requests
2, and an adaptive inner construction when it requests 4.  The scalar
adaptive counter requests exactly the outer side length many virtual 4s;
each costs exactly the number of cells of the inner box, giving one
physical 4 per cell overall.  Completed slices are invariant under every
inner move, while outer moves act on their coefficients exactly as the
scalar counter.  These two facts are the induction invariant.
\end{proof}

\begin{proof}[Completion of Theorem~\ref{thm:monotone}]
Lemma~\ref{lem:monoslices} constructs $D=4E$ with exactly $c$ fours, or
using only fours.  In either construction the first physical move merely
moves a single seed and spawns an equal second seed.  Delete that first
move and regard the resulting two equal tiles as the ordinary opening;
no score is lost.  This gives respectively a two-2 and two-4 opening.
The ledger identities yield the displayed score and length.  In the
all-4 construction the final mass is $4K$, so the length is $K-2$; no
ordinary opening has mass above 8 and no round adds more than 4 mass, so
$K-2$ is also a lower bound.  Finally, every line of $D$ is full with
unequal adjacent entries, so the endpoint is dead even if opposite
moves are subsequently allowed.
\end{proof}

\subsection{Reverse directions are genuinely necessary}

The monotone theorem also gives a clean obstruction to any purported
one-corner proof of the unrestricted cell-count ceiling.  Let $N_i^+$ be
the number of moves made in the direction opposite the chosen negative
direction on axis $i$.  Such a move can increase the lineage potential
$p$ by at most $n_i-1$.  Therefore every tile obeys
\[
 k+\sum_i x_i\le
 2+\sum_i(n_i-1)+\sum_i(n_i-1)N_i^+.
\]
Reflecting any subset of axes gives the analogous inequalities from every
corner.

\begin{corollary}[Direction obstruction]\label{cor:directions}
On a box with $d\ge2$ nontrivial axes, reaching the unrestricted
cell-count ceiling $2^{c+1}$ requires at least $d+1$ distinct direction
types during the play.  In particular two fixed directions can never
reach the $4\times4$ tile $131072$.
\end{corollary}
\begin{proof}
If an axis is never used, every lineage stays in a proper slice and is
subject to the smaller slice cell-count bound.  If exactly one direction
on each axis is used, Theorem~\ref{thm:monotone} bounds the exponent by
$2+\sum_i(n_i-1)<c+1$ for at least two nontrivial axes.
\end{proof}

For a $2\times n$ strip with Right never used, a target $2^{2n+1}$ at
cell $(r,c)$ therefore satisfies the sharper necessary inequalities
\[
 N_D\ge n-1+r+c,\qquad N_U\ge n-r+c,
\]
and hence $N_U+N_D\ge2n-1+2c$.  The finite three-direction strip
certificates described below therefore use the minimum possible number
of distinct direction types.

\subsection{A spawn-robust endgame on every rectangle}

Return now to unrestricted play.  Let $h\ge2$, $w\ge1$, and $c=hw$.
Define the full primed arrangement $Q_{h,w}$ by
\[
 Q(r,j)=2^{h(w-1-j)+r+1}\qquad(j<w-1),
\]
while its last column, from top to bottom, is
\[
 (4,4,8,16,\ldots,2^h).
\]
Its multiset is exactly $\{4,4,8,16,\ldots,2^c\}$ and its mass is
$2^{c+1}$.  For $4\times4$,
\[
 Q_{4,4}=\begin{pmatrix}
8192&512&32&4\\
16384&1024&64&4\\
32768&2048&128&8\\
65536&4096&256&16
\end{pmatrix}.
\]
For $h=2$ this is the canonical two-row ladder.

\begin{theorem}[Spawn-robust rectangular fold]\label{thm:rectfold}
From $Q_{h,w}$ the fixed word
\[
 (U^{h-1}L)^{w-1}U^{h-1}
\]
creates $2^{c+1}$ in exactly $c-1$ rounds for \emph{every} legal choice of
spawn cells and spawn values during those rounds.  No play from
$Q_{h,w}$ can create that tile sooner.
\end{theorem}
\begin{proof}
Designate one of the two 4s as the carry.  After $t$ useful carry merges,
its value is $2^{t+2}$.  All material spawned during the fold has total
mass at most $4t<2^{t+2}$ for $t\ge1$; initially there is none.  Maintain
a frontier column.  Columns to its left are full and unchanged.  In the
frontier, the carry is followed by its successive equal doubling partners;
all spawned material lies below this useful prefix or to its right.

An Up move merges the carry with its equal next partner.  Earlier columns
are full with unequal entries and therefore do not move.  Junk cannot
merge with or pass a useful tile because its entire mass is below the
current carry; the new spawn is necessarily outside the occupied useful
prefix.  After $h-1$ Up moves the useful material of the frontier column
has been consumed.  The next Left move merges the carry with the equal
top entry of the column immediately to the left, placing the doubled
carry at the new frontier while that column's remaining useful entries
stay below it.  The same invariant repeats one column left.  There are
$w(h-1)+(w-1)=c-1$ useful merges, giving the stated word.

For the lower bound, any target $2^{c+1}$ formed from the starting
position must contain either one of the initial 4s, or some subsequently
spawned 2/4: the other initial tiles have total mass only $2^{c+1}-8$.
A rank-2 leaf needs at least $c-1$ distinct rounds to reach rank $c+1$
(one doubling per round), and a rank-1 leaf needs even more.  Hence the
target cannot appear sooner.
\end{proof}

This theorem completely settles the endgame once the primed arrangement
is available.  It does \emph{not} prove that $Q_{h,w}$ is reachable from
an ordinary opening on every rectangle.  The unresolved part of the
Das--Paul induction is therefore a priming problem, not a folding problem.

\subsection{Near-ceiling rigidity in unrestricted play}

The proof of Theorem~\ref{thm:score} yields more than the numerical
ceiling.  Put
\[
 U_c=\Phi_c-4c=4(c-1)(2^c-1).
\]
At the $p$th tight-moment slot, the maximal rank is
$b=c+2-p$.  The maximizing score contribution is
$(b-1)2^b-4$; the best contribution after failing to saturate that slot is
$(b-2)2^{b-1}$.  Their difference is
\[
 b2^{b-1}-4.
\]

\begin{theorem}[Score rigidity]\label{thm:scorerigidity}
Suppose a legal $c$-cell play ends with score $U_c-\delta$.  If
$3\le a\le c+1$ and
\[
 \delta<a2^{a-1}-4,
\]
then its final board contains every tile
$2^a,2^{a+1},\ldots,2^{c+1}$.  In particular, a post-spawn board that is
neither of the two equality multisets in Theorem~\ref{thm:score} has
score at most $U_c-8$.
\end{theorem}
\begin{proof}
Every slot has nonnegative deficit from its independent maximum.  If a
slot of maximal rank $b\ge3$ is omitted or not saturated, it loses at
least $b2^{b-1}-4$.  This gap increases with $b$, so a total deficit below
the gap at $a$ forces every slot with $b\ge a$ to be saturated and hence
forces the listed ranks.  If all slots $p<c$ are saturated, the mandatory
last small spawn gives exactly one of the two equality multisets; otherwise
a slot with $b\ge3$ loses at least $3\cdot4-4=8$.
\end{proof}

\begin{corollary}[Length rigidity]\label{cor:lengthrigidity}
Let $L_c=2^{c+1}-c-4$.  If a legal play has length
$L_c-\varepsilon$ and
\[
 \varepsilon<2^{a-2}-1,
\]
then its final board contains every tile $2^a,\ldots,2^{c+1}$.  Outside
the two equality multisets, every play has length at most $L_c-1$.
\end{corollary}
\begin{proof}
The maximizing length contribution at a slot of maximal rank $b$ is
$2^{b-1}-1$ and the best nonsaturating contribution is $2^{b-2}$, a gap
of $2^{b-2}-1$.  Apply the same nonnegative-deficit argument.
\end{proof}

For $4\times4$, the known score $3{,}925{,}224$ has deficit 6876, while
$11\cdot2^{10}-4=11260$.  Thus every legal play meeting or beating that
score must end with all seven tiles
\[
131072,65536,32768,16384,8192,4096,2048.
\]
This is a compulsory structural feature of any improved search, not a
choice of heuristic target.

Let $q_F$ be the fewest spawned 4s among reachable full-chain games.  The
best score on that fixed chain is $\Phi_c-4q_F$, but this need not be the
global score optimum a priori.  Theorem~\ref{thm:scorerigidity} supplies a
sufficient reduction: if $q_F\le c+2$, the best full-chain score is a
global score optimum; if $q_F\le c+1$, the longest full-chain play is a
global length optimum.  Consequently an 18-four full-chain construction
on $4\times4$ would already justify the fixed-chain reduction for score,
while a 16-four construction would attain both universal ceilings.

\subsection{Finite strip frontier and executable checks}

The unrestricted two-row construction is not yet uniform, but explicit
independently replayed certificates now extend the finite frontier through
$2\times8$.  For $2\times n$, $2\le n\le8$, the cell-count maximum
$2^{2n+1}$ is reached in exactly
\[
 2^{2n-1}+2n-3
\]
rounds from a suitable two-4 opening, meeting Theorem~\ref{thm:moves}; a
suitable two-2 opening requires exactly one additional preparatory round.
The certificates use only Up, Left and Down, so Corollary~\ref{cor:directions}
shows that three distinct direction types are optimal.

\begin{center}\small
\begin{tabular}{@{}lrrr@{}}
\toprule
board & top tile & minimum rounds & two initial 2s\\
\midrule
$2\times2$ & 32 & 9 & 10\\
$2\times3$ & 128 & 35 & 36\\
$2\times4$ & 512 & 133 & 134\\
$2\times5$ & 2048 & 519 & 520\\
$2\times6$ & 8192 & 2057 & 2058\\
$2\times7$ & 32768 & 8203 & 8204\\
$2\times8$ & 131072 & 32781 & 32782\\
\bottomrule
\end{tabular}
\end{center}

The general monotone constructor, its standalone verifier, complete
negative-direction enumerations on several small boxes, and exhaustive
spawn-branch checks of Theorem~\ref{thm:rectfold} are under
\path{research/monotone/}.  These additions are executable checks of the
written proofs; they have not been translated into Lean.  In particular,
they do not expand the formal scope stated in Section~\ref{sec:lean}.

The unrestricted frontier remains explicit: the all-width priming
construction, the 16-four $4\times4$ full-chain game, and a stochastic
strong solution of the standard board remain open.  The results above
remove the monotone family and the entire rectangular endgame from that
list, and sharpen the structure any remaining construction must obey.
